% System definition and accounting; evaluation details are in Section V.
\section{System Model}\label{sec:system}

\subsection{Network and Query Model}
A UAV captures an aerial image, and an edge server answers a question
about the observed scene. The UAV selects either structured detection
evidence or an image for transmission. The edge server then applies
the corresponding decoder, as shown in Fig.~\ref{fig:arch}. The
objective is to answer the question with low communication and
computation energy, rather than to maximize image reconstruction quality.

Let $I$ denote the source image and $q$ the question. We consider five
structured question types: presence, counting, comparison, co-presence,
and threshold. The model assumes that the question and its type are
available before evidence selection. Query delivery and parsing are
outside the communication and computation budget considered here.
The type rule uses only the question type. Other selectors may also
use an assumed SNR bin or features extracted from the image, as
specified in Section~\ref{sec:routing}.

Each image--question pair is treated as a separate transaction, with
no amortization of detection or transmission across questions. The
evaluation uses aerial images and recorded inference outcomes; it
does not include video scheduling or a live UAV control loop.

\begin{figure*}[!t]\centering
\fbox{\parbox{0.89\linewidth}{\centering
Image and question $\longrightarrow$ sender-side evidence selection\\[5pt]
\begin{tabular}{ccl}
Detection records & $\longrightarrow$ wireless link $\longrightarrow$ & Symbolic decoder\\[3pt]
JPEG image & $\longrightarrow$ wireless link $\longrightarrow$ & Frozen VLM
\end{tabular}\\[5pt]
One selected payload and one receiver answering path per query}}
\caption{Evidence-level routing jointly selects the transmitted
representation and receiver computation. The type rule decides
before detection; the per-sample router first extracts detector
features. Neither policy requires an initial receiver VLM pass.}\label{fig:arch}
\end{figure*}

\subsection{Evidence Representation and Answer Generation}
\label{subsec:levels}\label{subsec:answering}
The detection-evidence branch uses an onboard
YOLOv8n detector~\cite{jocher2023yolov8} to extract
\begin{equation}
Z=g(I)=\{(c_j,b_j,p_j)\}_{j=1}^{J},
\end{equation}
where $J$ is the number of detected objects, and $c_j$, $b_j$, and
$p_j$ are the class label, bounding box, and confidence of object $j$.
These records form a symbolic packet. The receiver applies
question-specific count and logical operations to the received
evidence $\tilde Z$. No neural inference is required at the receiver
on this branch; neural detection is performed on the UAV.

The image branch transmits a JPEG representation of $I$. The
reconstructed image $\hat I$ and question $q$ are passed to a frozen
Qwen2-VL-2B model~\cite{wang2024qwen2vl}. The model generates the
answer by greedy decoding. Thus, evidence selection changes both the
transmitted representation and the computation used to answer the
question. Writing $s\in\mathcal{S}=\{\mathrm{det},\mathrm{img}\}$ for
the selected branch, the answer is
\begin{equation}\label{eq:answer}
\hat A_q=
\begin{cases}
D_{\mathrm{sym}}(\tilde Z,q), & s=\mathrm{det},\\[2pt]
D_{\mathrm{VLM}}(\hat I,q), & s=\mathrm{img},
\end{cases}
\end{equation}
where $D_{\mathrm{sym}}$ and $D_{\mathrm{VLM}}$ denote the symbolic
decoder and the VLM, respectively. Detector training, count
calibration, and answer scoring are described in
Section~\ref{subsec:setup}.

\subsection{Channel and Transmission Model}\label{subsec:phy}
\label{subsec:channels}
The uplink is represented by a normalized flat-fading channel,
\begin{equation}\label{eq:link}
y=hx+n,\qquad n\sim\mathcal{CN}(0,\sigma_n^2),
\end{equation}
with $\mathbb{E}[|x|^2]=1$ and $\mathbb{E}[|h|^2]=1$.
The average received SNR is $\gamma=1/\sigma_n^2$, and the
instantaneous SNR is $\gamma|h|^2$. We consider AWGN ($h=1$),
Rayleigh ($h\sim\mathcal{CN}(0,1)$), and Rician fading. For Rician
fading,
\begin{equation}
h=\sqrt{\frac{K}{K+1}}h_{\mathrm{LoS}}
 +\sqrt{\frac{1}{K+1}}h_{\mathrm{NLoS}},
\end{equation}
where $|h_{\mathrm{LoS}}|=1$,
$h_{\mathrm{NLoS}}\sim\mathcal{CN}(0,1)$, and $K$ is the linear
line-of-sight to scattered-power ratio. These models represent
different propagation conditions relevant to air-to-ground
links~\cite{khawaja2019survey}. Distance and path loss are absorbed
into the received SNR; UAV geometry and trajectory are not optimized.

\emph{Detection-evidence transmission.}\label{subsec:detphy}
Let $L_{\mathrm{det}}$ denote the packet length in bits. Its nominal
airtime uses rate-$1/2$ low-density parity-check (LDPC)
coding~\cite{gallager1962ldpc} and binary phase-shift keying (BPSK).
The main evaluation approximates evidence loss through
\begin{equation}\label{eq:outage}
P_{\mathrm{out}}(\gamma)
=\Pr\!\left[\log_2(1+\gamma|h|^2)
 <\frac{L_{\mathrm{det}}}{B\tau}\right],
\end{equation}
where $B$ is the bandwidth and $\tau$ is a reference transmission
window. This is an information-outage abstraction, not the measured
frame error rate (FER) of a compact LDPC packet. Its single fading
gain does not describe coding across multiple independent fading
blocks~\cite{tse2005fundamentals}.

The evidence impairment model allows records to be dropped or
corrupted, with equal probabilities conditional on a loss event.
Surviving records form $\tilde Z$ for symbolic decoding. This is an
application-level corruption model, not a model of residual errors
after packet integrity checks. The resulting answer accuracy is
conditional on this impairment assumption.

\emph{Image transmission.}\label{subsec:imgphy}
The rate-adaptive image branch uses the ergodic spectral efficiency
\begin{equation}\label{eq:se}
\mathrm{SE}(\gamma)
=\mathbb{E}_{h}[\log_2(1+\gamma|h|^2)]
\end{equation}
as an idealized rate for airtime accounting~\cite{tse2005fundamentals}.
The image traces use channel-dependent JPEG quality and resolution.
Let $L_{\mathrm{img}}(I,\gamma)$ be the compressed payload length
in bits before receiver-side decoding and file storage. The codec
reduces JPEG quality and, if needed, image resolution to meet the
reference transmission window. A received image saved again as a
JPEG is not used to estimate the transmitted payload.
The accounting uses the recorded reconstruction $\hat I$ for answer
evaluation and does not establish that a finite-length LDPC code
achieves \eqref{eq:se}. On an outage, the model returns a gray
image and charges the full reference window; no successful source
encoding length is recorded for that event.

\emph{Airtime.}\label{subsec:accounting}
With one complex channel use per $1/B$ seconds, the compact-packet
and ideal-rate accounting conventions give
\begin{equation}\label{eq:uses}
\begin{aligned}
T_s(\gamma)&=\frac{u_s(\gamma)}{B},\\
u_{\mathrm{det}}&=2L_{\mathrm{det}},\qquad
u_{\mathrm{img}}(\gamma)=
\frac{L_{\mathrm{img}}(I,\gamma)}{\mathrm{SE}(\gamma)}.
\end{aligned}
\end{equation}
Dependence on the source image is suppressed in $T_s$ and $u_s$ for
brevity. The image expression applies to non-outage transmissions; on outage,
$T_{\mathrm{img}}=\tau$. The energy tables use these conventions.
The nominal compact-packet airtime is not jointly coupled to the
outage-based reliability in \eqref{eq:outage}; a joint finite-length
reliability and airtime evaluation is needed to establish that
operating point. We therefore distinguish modeled airtime from a
validated link-level latency guarantee.

\subsection{Energy and Latency Models}\label{subsec:energyacct}
We account for communication and computation energy per question,
following the joint accounting approach of~\cite{dai2025pscom}.
Let $\rho$ denote the selector, and let $d_\rho(s)$ indicate whether
the detector runs for the selected branch. Type-based selectors run
the detector only when sending detection evidence. Per-sample
selectors need detector features before making a choice, so their
detector runs for every query. The platform-specific costs are
\begin{equation}\label{eq:energy}
\begin{aligned}
E_\rho^{\mathrm{UAV}}(s,\gamma)
 &=\frac{P_{\mathrm{tx}}}{\eta}T_s(\gamma)
   +d_\rho(s)E_{\mathrm{det}}\\
 &\quad+E_{\mathrm{route},\rho}+E_{\mathrm{enc}}(s),\\
E^{\mathrm{edge}}(s)
 &=\begin{cases}
 E_{\mathrm{sym}}, & s=\mathrm{det},\\
 E_{\mathrm{VLM}}, & s=\mathrm{img},
 \end{cases}\\
E_\rho(s,\gamma)
 &=E_\rho^{\mathrm{UAV}}(s,\gamma)+E^{\mathrm{edge}}(s).
\end{aligned}
\end{equation}
Here $P_{\mathrm{tx}}$ is radiated power, $\eta$ is power-amplifier
efficiency, and $E_{\mathrm{enc}}$ covers source encoding beyond
detector execution. $E_{\mathrm{det}}$, $E_{\mathrm{route},\rho}$,
$E_{\mathrm{sym}}$, and $E_{\mathrm{VLM}}$ denote detector, routing,
symbolic-decoding, and VLM-inference energy. The reference account
uses $\eta=1$. The efficiency term permits amplifier losses to be
included as in~\cite{cui2005energy}, but the reported comparisons
use the reference efficiency. Transmit power and the link budget
are not jointly optimized.

VLM energy is measured on an RTX~4060 edge-server proxy. Onboard
detector energy is modeled using a Jetson-class power and latency
range, with a separate desktop measurement as a cross-check.
Routing, symbolic decoding, and image source encoding are assigned
zero energy in the reported account; these are exclusions, not
measurements of zero consumption. Omitting image encoding reduces
the charged cost of the image branch. It does not bound all omitted
system costs. Numerical settings are given in
Section~\ref{subsec:energy-exp}.

For sequential execution, the corresponding per-question latency is
\begin{equation}\label{eq:latency}
\begin{aligned}
T_\rho^{\mathrm{total}}(s,\gamma)
 &=d_\rho(s)T_{\mathrm{det}}^{\mathrm{cmp}}
   +T_{\mathrm{route},\rho}^{\mathrm{cmp}}
   +T_{\mathrm{enc}}^{\mathrm{cmp}}(s)\\
 &\quad+T_s(\gamma)+T_{\mathrm{ans}}^{\mathrm{cmp}}(s),
\end{aligned}
\end{equation}
where the superscript $\mathrm{cmp}$ distinguishes processing time
from airtime. Answer processing is symbolic decoding or VLM
inference, according to the selected branch. Queueing and inter-query
overlap are not modeled. If transmission occupies the full reference
window, end-to-end latency also includes processing time and therefore
exceeds $\tau$.

The main account uses over-idle inference energy for an unbatched,
dedicated edge server.
Flight, image acquisition, control signaling, and radio circuitry
outside the amplifier model are excluded. Total accounted energy
includes both platforms, but UAV battery savings depend on the net
change in $E_\rho^{\mathrm{UAV}}$, including onboard computation.
Avoided edge inference does not directly extend UAV battery life.
